\documentclass[12pt,a4paper]{article}
\usepackage[utf8]{inputenc}
\usepackage[T1]{fontenc}
\usepackage[polish]{babel}
\usepackage{amsmath}
\usepackage{amsfonts}
\usepackage{amssymb}
\usepackage{booktabs}
\usepackage{geometry}
\usepackage{enumitem}
\usepackage{graphicx}
\usepackage{float}
\usepackage{listings}
\usepackage{xcolor}
\usepackage{tabularx}
\usepackage{array}
\usepackage{etoolbox}
\usepackage{url}

% lstlisting nigdy nie łamany między stronami
\BeforeBeginEnvironment{lstlisting}{\begin{minipage}{\linewidth}}
\AfterEndEnvironment{lstlisting}{\end{minipage}}

\geometry{margin=2cm}

\definecolor{codebg}{rgb}{0.97,0.97,0.97}
\definecolor{codekey}{rgb}{0.0,0.3,0.7}
\definecolor{codecomment}{rgb}{0.2,0.5,0.2}
\definecolor{codestring}{rgb}{0.6,0.1,0.1}

\lstset{
    language=C++,
    basicstyle=\ttfamily\small,
    keywordstyle=\color{codekey}\bfseries,
    commentstyle=\color{codecomment}\itshape,
    stringstyle=\color{codestring},
    backgroundcolor=\color{codebg},
    frame=single,
    framesep=4pt,
    breaklines=true,
    showstringspaces=false,
    numbers=left,
    numberstyle=\tiny\color{gray},
    xleftmargin=2em,
    tabsize=4,
    columns=fullflexible,
    literate=
        {ą}{{\k{a}}}1 {Ą}{{\k{A}}}1
        {ę}{{\k{e}}}1 {Ę}{{\k{E}}}1
        {ó}{{\'o}}1   {Ó}{{\'O}}1
        {ś}{{\'s}}1   {Ś}{{\'S}}1
        {ż}{{\.z}}1   {Ż}{{\.Z}}1
        {ź}{{\'z}}1   {Ź}{{\'Z}}1
        {ć}{{\'c}}1   {Ć}{{\'C}}1
        {ń}{{\'n}}1   {Ń}{{\'N}}1
        {ł}{{\l{}}}1  {Ł}{{\L{}}}1,
}

\title{Systemy czasu rzeczywistego, Projekt 5}
\author{Piotr Gugnowski, wydział EiTI, nr indeksu 320690}
\date{Czerwiec 2026}

\pagestyle{empty}

\begin{document}

\maketitle

% ===============================================================================
\section*{Zadanie 5.1: Standard IEEE~754: zakresy obu systemów kodowania}
% ===============================================================================

\subsection*{Treść zadania}

Wyznaczyć zakresy liczb w~obu systemach kodowania: kodzie
"wewnętrznym" oraz w~standardzie IEEE~754.

% -----------------------------------------------------------------------------
\subsection*{Analiza struktury bitowej}

Oba formaty są 32-bitowe. Struktura IEEE~754 single precision:
1~bit znaku, 8~bitów wykładnika z~nadmiarem~127, 23~bity mantysy
z~niejawnym bitem wiodącym.

Struktura kodu wewnętrznego \textbf{nie jest wprost podana w~treści
zadania}. Zadanie~5.1 wymaga jej jednoznacznego wyznaczenia przez
analizę wzorców bitowych z~tabeli konwersji. Wymaganiem jest,
aby format obsługiwał \textbf{liczby ujemne oraz wartości mniejsze od~1}
(ułamki), dokładnie tak jak IEEE~754.

\subsubsection*{Format IEEE~754 single precision}

\begin{center}
\begin{tabular}{|c|c|c|}
\hline
\textbf{Bit 31} & \textbf{Bity 30-23} & \textbf{Bity 22-0} \\
\hline
znak $S$ (1 bit) & wykładnik ze~stałą $E_b$ (8 bitów) & mantysa $M$ (23 bity) \\
\hline
\end{tabular}
\end{center}

Wartość liczby znormalizowanej:
\[
  (-1)^S \times (1 + M \cdot 2^{-23}) \times 2^{\,E_b - 127}
  \qquad E_b \in [1,\,254]
\]

Stała (bias) wykładnika wynosi $127$.

\subsubsection*{Kod "wewnętrzny"}

Na~podstawie tabeli konwersji (np.\ $1 \to \texttt{0x80000000}$,
$9 \to \texttt{0x83100000}$, $65535 \to \texttt{0x8F7FFF00}$) wyznaczono
następującą strukturę bitową:

\begin{center}
\begin{tabular}{|c|c|c|}
\hline
\textbf{Bity 31-24} & \textbf{Bit 23} & \textbf{Bity 22-0} \\
\hline
wykładnik ze~stałą $E_b$ (8 bitów, bias $128$) & znak $S$ (1 bit) & mantysa $M$ (23 bity) \\
\hline
\end{tabular}
\end{center}

Wartość liczby znormalizowanej:
\[
  (-1)^S \times (1 + M \cdot 2^{-23}) \times 2^{\,E_b - 128}
  \qquad E_b \in [1,\,254],\quad M \in [0,\,2^{23}-1]
\]

\paragraph{Wyznaczenie nadmiaru z~tabeli.}
Powyższą strukturę (w~tym wartość nadmiaru) odczytujemy wprost z~danych,
bez żadnych założeń. Po~pierwsze, mantysa (bity~22-0) jest w~obu kodach
identyczna, co widać np.\ dla liczby~9 (mantysa $=\texttt{0x100000}$ zarówno
w~\texttt{0x83100000}, jak i~\texttt{0x41100000}) oraz~65535
($\texttt{0x7FFF00}$ w~obu). Pozostaje wyznaczyć rolę najstarszych bitów.
Dla każdej liczby z~tabeli znamy jej rzeczywisty wykładnik~$E$ (z~postaci
znormalizowanej $1{,}M \times 2^E$). Porównajmy go z~8-bitowym polem w~bitach
[31:24] kodu wewnętrznego oraz, dla kontroli, z~polem wykładnika IEEE~754
(bity [30:23], o~znanym nadmiarze~$127$):

\begin{table}[H]
\centering
\small
\begin{tabular}{rrccc c}
\toprule
\textbf{Liczba} & $E$ & wewn.\ [31:24] & \textbf{wewn.\ $-\,E$} & IEEE [30:23] & IEEE $-\,E$ \\
\midrule
1     & 0  & $\texttt{0x80}=128$ & $\mathbf{128}$ & $\texttt{0x7F}=127$ & $127$ \\
2     & 1  & $\texttt{0x81}=129$ & $\mathbf{128}$ & $\texttt{0x80}=128$ & $127$ \\
9     & 3  & $\texttt{0x83}=131$ & $\mathbf{128}$ & $\texttt{0x82}=130$ & $127$ \\
65535 & 15 & $\texttt{0x8F}=143$ & $\mathbf{128}$ & $\texttt{0x8E}=142$ & $127$ \\
65536 & 16 & $\texttt{0x90}=144$ & $\mathbf{128}$ & $\texttt{0x8F}=143$ & $127$ \\
\bottomrule
\end{tabular}
\caption{Wyprowadzenie nadmiaru. Pole [31:24] kodu wewnętrznego pomniejszone
         o~rzeczywisty wykładnik daje \emph{stałą}~$128$ dla każdej liczby.}
\end{table}

Kolumna $\text{wewn.}-E$ jest stała i~równa~$128$ w~każdym wierszu. Dowodzi to,
że bity [31:24] tworzą 8-bitowe pole wykładnika ze~stałą (nadmiarem)
\textbf{128}. Analogiczna kolumna dla IEEE~754 jest stała i~równa~$127$
(znany nadmiar standardu). Różnica nadmiarów wynosi dokładnie~$1$ i~to właśnie
ona przekłada się dalej na relację wartości (połowa) i~prostą konwersję.
Bit~23 we~wszystkich przykładach jest zerem (liczby dodatnie); jako jedyny
wolny bit między wykładnikiem a~mantysą pełni rolę \textbf{bitu znaku},
analogicznie do bitu~31 w~IEEE~754, co umożliwia liczby ujemne.

\medskip
Format jest zatem \textbf{w~pełni analogiczny do IEEE~754}: definiuje te same
klasy liczb i~rezerwuje te same wartości pola wykładnika. Jedyne różnice to
wyznaczone wyżej dwa parametry: nadmiar ($128$ zamiast $127$) oraz położenie
bitu znaku (bit~23 zamiast bitu~31).

\begin{center}
\begin{tabular}{lll}
\toprule
$E_b$ & $M$ & Znaczenie (dokładnie jak w~IEEE~754) \\
\midrule
$0$        & $0$       & zero ze~znakiem: $+0 = \texttt{0x00000000}$, $-0 = \texttt{0x00800000}$ \\
$0$        & $\neq 0$  & liczba zdenormalizowana: $(-1)^S \times (M\cdot2^{-23}) \times 2^{-127}$ \\
$1\ldots254$ & dowolna & liczba znormalizowana: $(-1)^S \times (1 + M\cdot2^{-23}) \times 2^{E_b-128}$ \\
$255$      & $0$       & $\pm\infty$ \\
$255$      & $\neq 0$  & NaN \\
\bottomrule
\end{tabular}
\end{center}

\subsubsection*{Relacja wartości do IEEE~754 (klucz do konwersji)}

Ponieważ jedyną różnicą \emph{liczbową} jest nadmiar (większy o~$1$), wartość
słowa kodu wewnętrznego jest dokładnie \textbf{połową} wartości słowa IEEE~754
zbudowanego z~tych samych pól (znak, wykładnik, mantysa):
\[
  \text{wartość}_{\text{wew}}(X) = \tfrac{1}{2}\,\text{wartość}_{\text{IEEE}}\bigl(\text{te same pola}\bigr).
\]
Tożsamość ta zachodzi dla \emph{wszystkich} klas liczb (sprawdzenie:
dla znormalizowanej $2^{E_b-128} = 2^{E_b-127}/2$; dla zdenormalizowanej
$2^{-127} = 2^{-126}/2$; a~$\pm\infty/2 = \pm\infty$, $0/2 = 0$). Dzięki niej
konwersja w~obie strony sprowadza się do przeniesienia bitu znaku i~pomnożenia
(lub podzielenia) wartości przez~$2$, co automatycznie zachowuje semantykę
IEEE~754 dla znaku zera, liczb zdenormalizowanych, Inf/NaN, nadmiaru i~niedomiaru.

\subsubsection*{Przypadek specjalny: zero}

Zero jest przypadkiem szczególnym, bo nie da się go zapisać w~postaci
znormalizowanej $1{,}M \times 2^E$ (brak niejawnej jedynki wiodącej).
Tak jak w~IEEE~754, kod wewnętrzny rezerwuje na nie $E_b = 0$ (przy $M = 0$)
i~\textbf{rozróżnia znak zera}: $+0 = \texttt{0x00000000}$ oraz
$-0 = \texttt{0x00800000}$ (ustawiony bit znaku~23). Konwersja zachowuje znak
zera w~obie strony: $+0 \leftrightarrow \texttt{0x00000000}$,
$-0_{\text{IEEE}}\ (\texttt{0x80000000}) \leftrightarrow -0_{\text{wew}}\ (\texttt{0x00800000})$.
W~przykładach z~tabeli konwersji występuje wyłącznie $+0$.

Wszystkie liczby z~tabeli to dodatnie liczby całkowite, więc nie pokazują
jeszcze obsługi znaku ani ułamków, lecz wynika ona wprost ze~struktury:
\begin{itemize}
  \item \textbf{liczby ujemne}: bit znaku $S = 1$ (np.\ $-9 = \texttt{0x83900000}$),
  \item \textbf{ułamki $< 1$}: wykładnik ze~stałą $E_b < 128$, czyli $E < 0$
        (np.\ $0{,}5 = \texttt{0x7F000000}$).
\end{itemize}

% -----------------------------------------------------------------------------
\subsection*{Weryfikacja na przykładach z tabeli}

\begin{table}[H]
\centering
\small
\begin{tabular}{rlcccll}
\toprule
\textbf{Liczba} & \textbf{Kod wewn.} & $S$ & $E_b$ & $E$ & $M$ & \textbf{IEEE 754}\\
\midrule
0       & \texttt{0x00000000} & n/d & n/d & n/d & n/d & \texttt{0x00000000}\\
1       & \texttt{0x80000000} & 0  & 128 & 0   & 0                 & \texttt{0x3F800000}\\
2       & \texttt{0x81000000} & 0  & 129 & 1   & 0                 & \texttt{0x40000000}\\
9       & \texttt{0x83100000} & 0  & 131 & 3   & \texttt{0x100000} & \texttt{0x41100000}\\
65535   & \texttt{0x8F7FFF00} & 0  & 143 & 15  & \texttt{0x7FFF00} & \texttt{0x477FFF00}\\
65536   & \texttt{0x90000000} & 0  & 144 & 16  & 0                 & \texttt{0x47800000}\\
\midrule
\multicolumn{7}{l}{\textit{przykłady spoza tabeli, demonstracja ułamków i~liczb ujemnych:}}\\
$0{,}5$ & \texttt{0x7F000000} & 0  & 127 & $-1$ & 0                & \texttt{0x3F000000}\\
$-1$    & \texttt{0x80800000} & 1  & 128 & 0   & 0                 & \texttt{0xBF800000}\\
$-9$    & \texttt{0x83900000} & 1  & 131 & 3   & \texttt{0x100000} & \texttt{0xC1100000}\\
\bottomrule
\end{tabular}
\caption{Weryfikacja struktury bitowej. Dla każdej liczby
         $E_{b,\text{IEEE}} = E_b - 1$, a~mantysa $M$ jest identyczna w~obu formatach.}
\end{table}

\noindent Przykładowo dla $9 = 1{,}001_2 \times 2^3$:
\[
  S = 0,\quad E = 3 \Rightarrow E_b = 3 + 128 = 131 = \texttt{0x83},\quad
  M = 0{,}001_2 \Rightarrow M \cdot 2^{23} = 2^{20} = \texttt{0x100000}
\]
\[
  \text{wew} = (\texttt{0x83} \ll 24) \mid (0 \ll 23) \mid \texttt{0x100000}
  = \texttt{0x83100000} \checkmark
\]
\[
  E_{b,\text{IEEE}} = 131 - 1 = 130 = \texttt{0x82} \Rightarrow
  \text{IEEE} = (130 \ll 23) \mid \texttt{0x100000} = \texttt{0x41100000} \checkmark
\]

Dla liczby ujemnej $-9$ wystarczy ustawić bit znaku (bit~23):
$\text{wew} = \texttt{0x83100000} \mid \texttt{0x00800000} = \texttt{0x83900000}$.

% -----------------------------------------------------------------------------
\subsection*{Zakresy}

\subsubsection*{Kod "wewnętrzny"}

\begin{table}[H]
\centering
\begin{tabular}{ll}
\toprule
\textbf{Parametr} & \textbf{Wartość} \\
\midrule
Zakres & $\bigl[-(2-2^{-23})\cdot2^{126},\ (2-2^{-23})\cdot2^{126}\bigr]$ \\
Maksymalna wartość          & $(2 - 2^{-23}) \times 2^{126} \approx 1{,}701 \times 10^{38}$ \\
Min.\ wartość dodatnia znormalizowana & $2^{-127} \approx 5{,}877 \times 10^{-39}$ \\
Min.\ wartość dodatnia zdenormalizowana & $2^{-150} \approx 7{,}006 \times 10^{-46}$ \\
Zakres wykładnika $E$ (znorm.) & $[-127,\ 126]$ (8 bitów, bias $128$) \\
Zakres mantysy $M$          & $[0,\ 2^{23}-1]$ (23 bity) \\
Liczby ujemne / ułamki / Inf / NaN & tak (jak IEEE~754) \\
\bottomrule
\end{tabular}
\caption{Zakres kodu "wewnętrznego".}
\end{table}

\subsubsection*{Standard IEEE~754 single precision}

\begin{table}[H]
\centering
\begin{tabular}{ll}
\toprule
\textbf{Parametr} & \textbf{Wartość} \\
\midrule
Zakres & $\bigl[-(2-2^{-23})\cdot2^{127},\ (2-2^{-23})\cdot2^{127}\bigr]$ \\
Minimalna wartość dodatnia (znormalizowana) & $2^{-126} \approx 1{,}175 \times 10^{-38}$ \\
Maksymalna wartość & $(2 - 2^{-23}) \times 2^{127} \approx 3{,}403 \times 10^{38}$ \\
Zakres wykładnika & $E \in [-126,\ +127]$ (bias $= 127$) \\
Zakres mantysy $M$ & $[0,\ 2^{23}-1]$ (23 bity) \\
Liczby ujemne & tak (bit znaku, bit~31) \\
\bottomrule
\end{tabular}
\caption{Zakres standardu IEEE~754 single precision (32-bit).}
\end{table}

\subsubsection*{Porównanie zakresów}

\begin{table}[H]
\centering
\begin{tabular}{lll}
\toprule
\textbf{Cecha} & \textbf{Kod wewnętrzny} & \textbf{IEEE 754} \\
\midrule
Liczby ujemne   & tak (bit~23) & tak (bit~31) \\
Wartości $< 1$  & tak & tak \\
Zero ze~znakiem ($\pm0$) & tak & tak \\
Liczby zdenormalizowane & tak & tak \\
Inf / NaN & tak & tak \\
Bity wykładnika & 8 (bias $128$) & 8 (bias $127$) \\
Bity mantysy    & 23 & 23 \\
Zakres wykładnika $E$ (znorm.) & $[-127,\ 126]$ & $[-126,\ 127]$ \\
Maksymalna wartość & $\approx 1{,}70 \times 10^{38}$ & $\approx 3{,}40 \times 10^{38}$ \\
\bottomrule
\end{tabular}
\caption{Porównanie cech obu systemów kodowania.}
\end{table}

Oba formaty są \textbf{identyczne semantycznie}: te same klasy liczb (zero ze~znakiem,
zdenormalizowane, znormalizowane, Inf/NaN), te same zarezerwowane wartości
wykładnika, to samo 23-bitowe pole mantysy. Jedyna różnica liczbowa to nadmiar:
bias~$128$ przesuwa cały zakres wykładnika o~jedną binadę w~dół względem IEEE~754
($E \in [-127, 126]$ zamiast $[-126, 127]$). W~praktyce kod wewnętrzny sięga
nieco niżej przy małych wartościach, a~IEEE~754 nieco wyżej przy dużych
(maksimum większe o~czynnik~$2$). Jest to bezpośrednia, nieunikniona konsekwencja
nadmiaru~$128$ wynikającego z~tabeli konwersji w~treści zadania.

% ===============================================================================
\section*{Zadanie 5.2: Konwersja kodu wewnętrznego do IEEE~754}
% ===============================================================================

\subsection*{Treść zadania}

Napisać aplikację konwersji kodu "wewnętrznego" liczb
zmiennoprzecinkowych do~kodu liczb zmiennoprzecinkowych zgodnych z~IEEE~754.
Aplikacja ma umożliwiać wprowadzenie z~klawiatury liczb w~kodzie binarnym
naturalnym.

\subsection*{Implementacja}

Aplikacja napisana w~C++17, kompilowalna na systemach Windows, macOS i~Linux
(clang++ lub g++). Do budowania wymagany jest \texttt{just} (\texttt{just build}).
Kod źródłowy składa się z~dwóch plików:

\begin{itemize}
  \item \texttt{src/shared.cpp} / \texttt{src/shared.h}: funkcje wspólne
        (konwersje, wyświetlanie, parsowanie wejścia),
  \item \texttt{src/task2\_internal\_to\_ieee.cpp}: punkt wejścia (\texttt{main}).
\end{itemize}

\subsubsection*{Format wejścia}

Zgodnie z~założeniem zadania użytkownik wprowadza kod wewnętrzny jako
\textbf{32-bitowy ciąg w~kodzie binarnym naturalnym} (ciąg znaków \texttt{0}/\texttt{1}).
Spacje są dozwolone i~ignorowane, co umożliwia grupowanie bitów dla czytelności.

Przykład dla liczby~9, której kod wewnętrzny to \texttt{0x83100000}:


\begin{lstlisting}[language={},numbers=none]
$ ./task2_internal_to_ieee 10000011000100000000000000000000
$ ./task2_internal_to_ieee "1000 0011 0001 0000 0000 0000 0000 0000"
\end{lstlisting}

Aplikacja obsługuje dwa tryby:
\begin{enumerate}[label=\arabic*.]
  \item \textbf{Argument wiersza poleceń}: 32-bitowy ciąg binarny podany bezpośrednio.
  \item \textbf{Tryb interaktywny}: aplikacja wyświetla prompt i~czeka na wejście.
\end{enumerate}

\subsubsection*{Algorytm konwersji: kod wewnętrzny $\to$ IEEE~754}

Wykorzystujemy tożsamość $\text{wartość}_{\text{wew}} = \tfrac{1}{2}\,\text{wartość}_{\text{IEEE}}$
(z~sekcji~5.1). Algorytm:
\begin{enumerate}[label=\arabic*.]
  \item Przełóż pola słowa wewnętrznego do~układu IEEE: znak z~bitu~23 na~bit~31,
        wykładnik z~bitów~[31:24] na~[30:23], mantysa bez zmian.
  \item Pomnóż tak uzyskaną wartość przez~$\tfrac{1}{2}$ (zmiana nadmiaru
        $128 \to 127$). Wynikiem jest słowo IEEE~754.
\end{enumerate}

Dla liczb znormalizowanych krok~2 to po prostu zmniejszenie pola wykładnika
o~$1$ ($E_{b,\text{IEEE}} = E_b - 1$) przy niezmienionej mantysie. Mnożenie
realizowane sprzętowo (typ \texttt{float}) jest dokładne i~automatycznie
odtwarza zachowanie IEEE~754 dla wszystkich pozostałych klas: znaku zera,
liczb zdenormalizowanych, Inf/NaN oraz ewentualnego niedomiaru.

\begin{lstlisting}
uint32_t internal_to_ieee754(uint32_t internal) noexcept {
    const uint32_t S    = (internal & 0x00800000u) ? 1u : 0u; // bit 23
    const uint32_t E_b  = (internal >> 24u) & 0xFFu;          // bity [31:24]
    const uint32_t mant = internal & 0x007FFFFFu;             // bity [22:0]
    // te same pola w ukladzie IEEE (znak 31, wykladnik 30:23, mantysa 22:0)
    const uint32_t as_ieee = (S << 31u) | (E_b << 23u) | mant;
    // wartosc_wew = wartosc_ieee / 2   (nadmiar 128 -> 127)
    return float_to_bits(bits_to_float(as_ieee) * 0.5f);
}
\end{lstlisting}

\subsubsection*{Wyniki dla danych wzorcowych i~nowych przypadków}

\begin{table}[H]
\centering
\small
\begin{tabular}{rlll}
\toprule
\textbf{Liczba} & \textbf{Wejście (kod binarny naturalny)} & \textbf{Kod wewn. (hex)} & \textbf{IEEE~754 (hex)} \\
\midrule
0     & \texttt{0000...0} (32 zera)              & \texttt{0x00000000} & \texttt{0x00000000} \\
1     & \texttt{1000 0000 0000 ... 0000}         & \texttt{0x80000000} & \texttt{0x3F800000} \\
9     & \texttt{1000 0011 0001 0000} \ldots      & \texttt{0x83100000} & \texttt{0x41100000} \\
65535 & \texttt{1000 1111 0111 1111} \ldots      & \texttt{0x8F7FFF00} & \texttt{0x477FFF00} \\
65536 & \texttt{1001 0000 0000 ... 0000}         & \texttt{0x90000000} & \texttt{0x47800000} \\
\midrule
$0{,}5$ & \texttt{0111 1111 0000 ... 0000}       & \texttt{0x7F000000} & \texttt{0x3F000000} \\
$-1$    & \texttt{1000 0000 1000 ... 0000}       & \texttt{0x80800000} & \texttt{0xBF800000} \\
$-9$    & \texttt{1000 0011 1001 0000} \ldots    & \texttt{0x83900000} & \texttt{0xC1100000} \\
\bottomrule
\end{tabular}
\caption{Wyniki konwersji: wejście w~kodzie binarnym naturalnym, wyjście IEEE~754.
         Ostatnie trzy wiersze demonstrują obsługę ułamków i~liczb ujemnych.}
\end{table}

% ===============================================================================
\section*{Zadanie 5.3: Konwersja IEEE~754 do kodu wewnętrznego}
% ===============================================================================

\subsection*{Treść zadania}

Napisać aplikację konwersji liczb zmiennoprzecinkowych zgodnych z~IEEE~754
do~postaci kodu "wewnętrznego".

\subsection*{Implementacja}

Aplikacja napisana w~C++17, kompilowalna na systemach Windows, macOS i~Linux
(clang++ lub g++). Kod źródłowy:

\begin{itemize}
  \item \texttt{src/shared.cpp} / \texttt{src/shared.h}: funkcje wspólne
        (te same co w~zadaniu~5.2),
  \item \texttt{src/task3\_ieee\_to\_internal.cpp}: punkt wejścia (\texttt{main}).
\end{itemize}

\subsubsection*{Format wejścia}

Użytkownik wprowadza \textbf{32-bitową wartość IEEE~754 w~postaci szesnastkowej}
(z~przedrostkiem \texttt{0x} lub bez). Aplikacja odczytuje wzorzec bitowy słowa
i~konwertuje go na kod wewnętrzny, wypisując wynik również w~postaci szesnastkowej.

Przykłady wywołania:


\begin{lstlisting}[language={},numbers=none]
$ ./task3_ieee_to_internal 0x41100000     # 9
$ ./task3_ieee_to_internal 0x3F000000     # 0,5
$ ./task3_ieee_to_internal 0xC1100000     # -9
\end{lstlisting}

Aplikacja obsługuje dwa tryby:
\begin{enumerate}[label=\arabic*.]
  \item \textbf{Argument wiersza poleceń}: wartość szesnastkowa podana bezpośrednio.
  \item \textbf{Tryb interaktywny}: aplikacja wyświetla prompt i~czeka na wejście.
\end{enumerate}

\subsubsection*{Zachowanie dla wszystkich klas liczb}

Konwersja jest \emph{totalna}: każda wartość IEEE~754 ma odpowiednik w~kodzie
wewnętrznym, a~zachowanie dla każdej klasy liczb jest \textbf{takie samo jak
w~IEEE~754}:

\begin{center}
\begin{tabular}{ll}
\toprule
\textbf{Wejście IEEE~754} & \textbf{Wynik (kod wewnętrzny)} \\
\midrule
$\pm0$, liczba znormalizowana, zdenormalizowana & ta sama wartość (znak zachowany) \\
$\pm\infty$, NaN                                & $\pm\infty$, NaN \\
$|x| \geq 2^{127}$ (powyżej zakresu znorm.\ wewn.) & nadmiar $\to \pm\infty$ (jak w~IEEE) \\
\bottomrule
\end{tabular}
\end{center}

Nadmiar do~$\pm\infty$ przy największych liczbach IEEE~754 ($E = 127$) to
dokładnie takie samo zachowanie, jakie standard przewiduje przy zawężającej
konwersji (np.\ \texttt{double}~$\to$~\texttt{float}). Nie wprowadzamy żadnej
własnej reguły.

Niedomiar w~tym kierunku \emph{nie występuje}: zakres kodu wewnętrznego sięga
niżej niż IEEE~754 (do~$2^{-150}$ wobec $2^{-149}$), więc każda liczba IEEE~754,
łącznie ze~wszystkimi zdenormalizowanymi, jest reprezentowalna. Niedomiar
pojawia się dopiero w~kierunku odwrotnym (zadanie~5.2), gdy najmniejsze liczby
zdenormalizowane kodu wewnętrznego po~podzieleniu przez~$2$ spadają poniżej
najmniejszej liczby IEEE~754.

\subsubsection*{Algorytm konwersji: IEEE~754 $\to$ kod wewnętrzny}

Symetrycznie do zadania~5.2 korzystamy z~tożsamości
$\text{wartość}_{\text{IEEE}} = 2 \cdot \text{wartość}_{\text{wew}}$:
mnożymy wartość przez~$2$ (nadmiar $127 \to 128$), a~następnie przekładamy pola
do~układu wewnętrznego (znak z~bitu~31 na~bit~23). Dla liczb znormalizowanych
sprowadza się to do zwiększenia pola wykładnika o~$1$ przy niezmienionej mantysie.

\begin{lstlisting}
uint32_t ieee754_to_internal(uint32_t ieee) noexcept {
    // wartosc_ieee(as_ieee) = 2 * wartosc_wejscia   (nadmiar 127 -> 128)
    const uint32_t as_ieee = float_to_bits(bits_to_float(ieee) * 2.0f);
    const uint32_t S    = (as_ieee >> 31u) & 1u;        // bit 31
    const uint32_t E_b  = (as_ieee >> 23u) & 0xFFu;     // bity [30:23]
    const uint32_t mant = as_ieee & 0x007FFFFFu;        // bity [22:0]
    // uklad wewnetrzny: wykladnik w [31:24], znak na bit 23
    return (E_b << 24u) | (S << 23u) | mant;
}
\end{lstlisting}

Mnożenie i~dzielenie przez~$2$ realizowane sprzętowo (typ \texttt{float}) są
dokładne (zmieniają jedynie pole wykładnika) i~przenoszą całą semantykę IEEE~754
na~kod wewnętrzny bez dodatkowego kodu obsługi przypadków szczególnych.

\subsubsection*{Wyniki dla danych wzorcowych i~nowych przypadków}

\begin{table}[H]
\centering
\begin{tabular}{rll}
\toprule
\textbf{Liczba} & \textbf{IEEE~754 (wejście)} & \textbf{Kod wewnętrzny (wyjście)} \\
\midrule
0       & \texttt{0x00000000} & \texttt{0x00000000} \\
1       & \texttt{0x3F800000} & \texttt{0x80000000} \\
2       & \texttt{0x40000000} & \texttt{0x81000000} \\
9       & \texttt{0x41100000} & \texttt{0x83100000} \\
65535   & \texttt{0x477FFF00} & \texttt{0x8F7FFF00} \\
65536   & \texttt{0x47800000} & \texttt{0x90000000} \\
\midrule
$0{,}5$  & \texttt{0x3F000000} & \texttt{0x7F000000} \\
$0{,}25$ & \texttt{0x3E800000} & \texttt{0x7E000000} \\
$-1$     & \texttt{0xBF800000} & \texttt{0x80800000} \\
$-9$     & \texttt{0xC1100000} & \texttt{0x83900000} \\
$\pi$    & \texttt{0x40490FDB} & \texttt{0x81490FDB} \\
\bottomrule
\end{tabular}
\caption{Wyniki konwersji IEEE~754 $\to$ kod wewnętrzny. Ostatnie pięć wierszy
         demonstruje obsługę ułamków i~liczb ujemnych.}
\end{table}

\end{document}
